\section{Related Work}

% \gf{Need to discuss data distillation methods in this section.}
% \begin{enumerate}
%     \item Federated Learning related work
%     \item Differential Privacy related work, Zinan's paper
%     \item Synthetic data helping LLMs (Alpaca, Unnatural code, etc etc) 
%     \item Large language models
% \end{enumerate}
% We focus on a setting where the data is distributed among many client devices which do not want to leak sensitive/private information.  This has connections to the following fields.

\textbf{DP Federated Learning.} Differentially Private Federated Learning (FL) is a widely-used approach for learning ML models from distributed private data \citep{mcmahan2017learning,kairouz2021advances}. In DP-FL, model weights are sent to users who train the model locally on their private data; private local model updates are then collected at a central server. 
% A central topic in the area of FL involves understanding algorithms and tradeoffs between model quality and cost (i.e., client computation and communication), particularly in the face of privacy constraints \citep{hou2021fedchain, wang2020tackling, li2020federated, karimireddy2020scaffold, wang2021field,mishchenko2022proxskip,sadiev2022communication,grudzien2023can}. 
% For example, \citet{mcmahan2017learning}  demonstrated that DP FL optimization algorithms can be used to produce high-quality federated models with strong user-level privacy guarantees.
% Since then, 
Researchers have many techniques for improving privacy-utility tradeoffs in DP-FL. Some include shuffle-based privacy amplification  \cite{secagg,girgis2021shuffled,agarwal2018cpsgd}, pretraining on public datasets \cite{xu2023federated}, private selection of the best pretraining datasets \cite{hou2023privately,gu2022choosing}, and DP-FL methods that do not rely on uniform sampling/shuffling \citep{kairouz2021practical}.

\begin{figure}
    \centering
    \includegraphics[width=0.85\columnwidth]{figures/num_samples_ablation.pdf}
    \caption{Next-token prediction accuracy for PrE-Text as we vary the number of synthetic examples generated by the \texttt{Expand} part of the algorithm. We find that increasing the number of synthetic examples across several orders of magnitude improves the accuracy of the downstream model (DistilGPT2) roughly log-linearly, though the growth peaks on \textsc{Code} after 1M samples.
    }
    \label{fig: scaling}
\end{figure}

Today, growing efforts study how to train larger models on client data.  \citet{charles2022federated} propose to have users optimize slices of large models, though they have not demonstrated the approach on models larger than shallow logistic regression and convolutional neural network models.  \citet{collins2023profit, cai2022fedadapter, zhang2023towards, zhao2022reduce, guo2023promptfl} only tune sub-components of the models in the federated setting to reduce client computational burden and communication, but this still requires having clients store and perform inference with large models on their device.  \citet{zhang2023gpt} use foundation models to produce synthetic data to pretrain smaller models.  None of these methods consider privacy.
% \gf{To our knowledge, prior work does not provide a clear path for training large models over federated client data, even without privacy.}
% \gf{Maybe we can cut the above sentence. It's useful, but not a must-have imo}


% There are three main lines of work in this area: improving algorithmic efficiency \citep{woodworth2020local, karimireddy2020scaffold, hou2021fedchain}, personalization \citep{fallah2020personalized, jiang2019improving} and privacy \gf{cite}.  \gf{need more citations}

% \citet{mcmahan2017learning} demonstrate that one can combine FL with differential privacy (DP) to produce models that have privacy guarantees for users.  \citet{secagg} introduce secure aggregation, which allows FL algorithms that utilize averaging (like FedAvg \citep{pmlr-v54-mcmahan17a}) to amplify their privacy protections by using cryptographic secure multi-party computations.  \citep{xu2023federated} show that pretraining models on public data before FL training can improve privacy-utility trade-offs and \citep{hou2023privately} show how to choose the best pretraining dataset for a particular FL application in a differentially private way.  \citet{augenstein2019generative} generated DP synthetic image data by training GANs \citep{goodfellow2014generative} in a federated manner. 

\textbf{Synthetic Data.} Synthetic data is increasingly being used to train language models \citep{alpaca, wang2022self, honovich2022unnatural}.  Common approaches involve carefully designing prompts to ChatGPT \citep{radford2019language} to generate synthetic training data for another open source model, e.g. LLaMA \citep{touvron2023llama}, to replicate ChatGPT behavior.
In these works, the synthetic data is used solely to enhance final model utility, and does not satisfy any formal privacy guarantees. In the image setting, useful synthetic data is often produced using dataset distillation \citep{wang2018dataset, zhao2020dataset, cazenavette2022dataset}. This approach has been adapted to the federated setting \citep{song2023federated}.


\textbf{DP Synthetic Data.} Much of the work on producing \emph{private} (i.e., DP) synthetic data from deep generative models is in the image domain \citep{lin2021privacy, cao2021don,dockhorn2022differentially,chen2022fedtune}. 
Common approaches involve training a generative adversarial network (GAN) or diffusion model in a differentially private way, e.g., via DP-SGD, a DP version of stochastic gradient descent (SGD) \citep{abadi2016deep}.  \citet{tang2023privacy} generate DP synthetic text in the federated setting, but their method requires users to send private data to ChatGPT (which is located on-server).  
Such actions are not allowed under our threat model, where the central server is not trusted to hold private data. 
In concurrent work, \citep{xie2024differentially} propose a  private evolution algorithm for DP synthetic text data. Their method is similar to ours (we summarize differences in \cref{sec:intro}); the most important difference is that our focus is on understanding the relation between synthetic data and on device training, while their work aims to improve the quality of DP synthetic text data more generally.
Also in the text domain, \citep{li2021large, yu2021differentially} demonstrate that it is possible to finetune pretrained large language models in a central-DP manner (i.e., the private data is available to the model developer).

Recently, %some concurrent work released after the submission of our paper to ICML has 
several papers (some of them concurrent to ours) have proposed to finetune LLMs on DP synthetic data in the central DP setting \citep{yue2022synthetic, kurakin2023harnessing,yu2024privacy, ding2024delving}. 
In the private federated setting, work by \citet{wang2023can, wu2024prompt}
 also considers using LLM synthetic data. They propose to filter LLM synthetic data to samples that are relevant to the private data. The model is subsequently finetuned using DP-FL on-device. 
Their work shows a substantial improvement in next-word prediction accuracy compared to vanilla pre-training with DP-FL finetuning. In contrast, our work uses DP synthetic data to replace the DP-FL step itself, while keeping pre-training untouched. These methods are complementary and could potentially be combined.
Overall, the literature suggests that DP synthetic data may be useful for training models from private, client data.
% Given that they already assume users are uploading private data to the server, their setting 
% This is perhaps closer to the central DP setting (where the central server is completely trusted) than the federated DP setting (where the central server is considered to be honest but not trusted enough to hold private data).  

% so it is unknown how well their synthetic data will serve as training data.

% \textbf{DP Language Models.} In the text domain, \citep{li2021large, yu2021differentially} demonstrate that it is possible to finetune pretrained large language models in a central-DP manner (i.e., the private data is available to the model developer). This is a different privacy setting from ours. 
